\chapter{Feb.~24 --- Differentials}

\section{Review of Tangent Spaces}

\begin{definition}
  For a variety $X$ and $x \in X$,
  the \emph{tangent space} of $X$ at $x$ is
  \[
    T_{X, x}
    = (\m_x / \m_x^2)^\vee.
  \]
  We say that $X$ is \emph{smooth} at
  $x$ if $\dim_x X = \dim_k T_{X, x}$.
\end{definition}

\begin{remark}
  Our goals will be:
  \begin{enumerate}
    \item Define $\Omega_X \in \Coh(X)$
      such that $\Omega_X(U)$ gives the
      ``algebraic $1$-forms on $U$.''
    \item Define $T_X := \Omega_X^\vee$.
    \item Show that there is a natural
      isomorphism $T_X(x) \cong T_{X, x}$.
    \item Construct a relative version
      $\Omega_{X / Y}$ for a morphism
      $X \to Y$.
  \end{enumerate}
\end{remark}

\section{Differentials}

\begin{remark}
  Let $\phi : A \to B$ be a morphism
  of rings, e.g.
  $k \to k[x_1, \dots, x_n] / I$.
\end{remark}

\begin{definition}
  For a $B$-module $M$, an
  \emph{$A$-derivation} is a map
  $D : B \to M$ such that
  \begin{enumerate}
    \item $D$ is $A$-linear,
    \item (Leibniz rule)
      $D(b_1 b_2) = b_1 D(b_2) + D(b_1) b_2$.
  \end{enumerate}
\end{definition}

\begin{remark}
  Using (2), condition
  (1) can be reformulated as
  the following:
  \begin{enumerate}
    \item[1'.] $D$ is additive and
      $D(a \cdot 1) = 0$ for all $a \in A$.
  \end{enumerate}
\end{remark}

\begin{example}
  Let $A = k$ and $B = k[x_1, \dots, x_n]$.
  We can define
  \begin{align*}
    D : k[x_1, \dots, x_n]
    &\longrightarrow \bigoplus_{i = 1}^n k[x_1, \dots, x_n] dx_i \\
    f &\longmapsto df = \sum \frac{\partial f}{\partial x_i} dx_i.
  \end{align*}
  One can check that $D$ is a $k$-derivation
  (but $D$ is not $k[x_1, \dots, x_n]$-linear).
\end{example}

\begin{prop}
  There exists a derivation
  $B \to \Omega_{B / A}$ such that for
  any derivation $D : B \to M$, there exists
  unique a $B$-module homomorphism such that
  the following diagram commutes:
  \[
    \begin{tikzcd}
      B \ar[r, "d"] \ar[dr, "D", swap]
      & \Omega_{B / A} \ar[d, dashed, "\exists !"] \\
      & M
    \end{tikzcd}
  \]
  Here $\Omega_{B / A}$ is called the
  \emph{module of K\"ahler differentials}.
\end{prop}

\begin{proof}
  Let $\Omega_{B / A}$ be the $B$-module
  with generators $(db)_{b \in B}$
  and relations
  \begin{enumerate}
    \item $d(b_1 b_2) = b_1 d(b_2) + b_2 d(b_1)$ for $b_1, b_2 \in B$,
    \item $d(b_1 + b_2) = db_1 + db_2$
      for $b_1, b_2 \in B$
      and $d(\phi(a)) = 0$ for $a \in A$.
  \end{enumerate}
  Then one can check
  $B \to \Omega_{B / A}$ given by
  $b \mapsto db$ is a derivation
  satisfying the universal property.
\end{proof}

\begin{remark}
  If $(b_i)_{i \in I}$ generate $B$
  as an $A$-module, then the $db_i$
  generate $\Omega_{B / A}$ as a $B$-module.
\end{remark}

\begin{example}
  We have the following:
  \begin{enumerate}
    \item[0.] $\Omega_{A / A} = 0$.
    \item For $K \subseteq L$ a finite
      field extension, $\Omega_{L / K} = 0$
      if and only if $L / K$ is separable.
    \item Let $B = A[x_1, \dots, x_n]$.
      We know that $dx_i$ for
      $1 \le i \le n$ generate
      $\Omega_{B / A}$. We claim that the
      $dx_i$ are also linearly
      independent over $B$.

      To see this, fix $1 \le i \le n$ and
      consider the
      derivation $D_i : B \to B$ given by
      $f \mapsto \partial f / \partial x_i$.
      Then
      \[
        D_i(x_j) =
        \begin{cases}
          1 & \text{if $j = i$} \\
          0 & \text{otherwise}.
        \end{cases}
      \]
      By the universal property we have a
      diagram
      \[
        \begin{tikzcd}
          & \Omega_{B / A} \ar[d, dashed] \\
          B \ar[ur] \ar[r, "D_i", swap] & B
        \end{tikzcd}
      \]
      so the $dx_i$ are independent over
      $B$. Thus we have
      $\Omega_{B / A} = B dx_1 \oplus \cdots \oplus B dx_n$.
  \end{enumerate}
\end{example}

\begin{prop}
  If $S \subseteq B$ is a multiplicative
  system, then there is a canonical isomorphism
  \begin{align*}
    \Omega_{S^{-1} B / A}
    &\overset{\cong}{\longrightarrow} S^{-1} \Omega_{B / A} \\
    d(b / s) &\longmapsto \frac{1}{s} db - \frac{b}{s^2} ds.
  \end{align*}
\end{prop}

\begin{proof}
  Given an $A$-derivation
  $S^{-1} B \to M$, we get an $A$-derivation
  by composition:
  \[
    B \longrightarrow S^{-1} B \longrightarrow M
  \]
  So there exists
  a unique $B$-module homomorphism
  $D$ such that the following
  diagram commutes:
  \pagebreak
  \[
    \begin{tikzcd}
      B \ar[d] \ar[r] & \Omega_{B / A} \ar[d, "D", dashed] \ar[dr] \\
      S^{-1} B \ar[r] \ar[rr, bend right, "S^{-1} D", swap] & M & S^{-1} \Omega_{B / A} \ar[l, dashed, "\phi"]
    \end{tikzcd}
  \]
  Note that the universal property of
  localization,
  we get a unique map $S^{-1} D$ that
  fits into the above diagram. Then
  one can check that the map
  \begin{align*}
    S^{-1} D : S^{-1} B &\longrightarrow S^{-1} \Omega_{B / A} \\
    \frac{b}{s} &\longmapsto \frac{1}{s} db - \frac{b}{s^2} ds
  \end{align*}
  makes the diagram commute and is
  an $A$-derivation, so
  $S^{-1} \Omega_{B / A} \cong \Omega_{S^{-1} B / A}$.
\end{proof}

\begin{prop}\label{prop:properties-differentials}
  We have the following:
  \begin{enumerate}
    \item For $A$-algebras $B_1$ and
      $B_2$, we have
      $\Omega_{B_1 \otimes_A B_2} \cong \Omega_{B_1 / A} \otimes_A \Omega_{B_2 / A}$.
    \item If we have ring homomorphisms
      $A \overset{\phi}{\longrightarrow} B \overset{\psi}{\longrightarrow} C$, then
      the sequence
      \[
        \begin{tikzcd}
          \Omega_{B / A} \otimes_B C \ar[r]
          & \Omega_{C / A} \ar[r] &
          \Omega_{C / B} \ar[r] & 0
        \end{tikzcd}
      \]
      is exact, where the
      first map $\Omega_{B / A} \otimes_B C \to \Omega_{C / A}$ is given by
      $db \otimes c \mapsto c d\psi(b)$.
    \item If $I \le B$ is an ideal, then
      there exists an exact sequence
      \[
        \begin{tikzcd}
          I / I^2 \ar[r]
          & \Omega_{B / A} \otimes (B / I) \ar[r]
          & \Omega_{(B / I) / A} \ar[r]
          & 0
        \end{tikzcd}
      \]
      where the first map
      $I / I^2 \to \Omega_{B / A} \otimes (B / I)$ is given by
      $\overline{b} \mapsto db \otimes 1$.
  \end{enumerate}
\end{prop}

\begin{proof}
  See Mustata's notes.
\end{proof}

\begin{example}\label{example:affine-variety-differentials}
  Let $X \subseteq \Affine^n$ be an affine
  variety, $C = \OO_X(X) = k[x_1, \dots, x_n] / I(X)$,
  $B = k[x_1, \dots, x_n]$, and $A = k$.
  Let $I = I(X)$. Then
  Proposition \ref{prop:properties-differentials}(3) gives an exact sequence
  \[
    \begin{tikzcd}
      I / I^2 \ar[r]
      & \bigoplus_{i = 1}^n Bdx_i \otimes_B C \ar[r]
      & \Omega_{C / k} \ar[r] & 0
    \end{tikzcd}
  \]
  Thus we see that $\Omega_{C / k}$ is
  given by
  \[
    \Omega_{C / k}
    = \frac{\bigoplus_{i = 1}^n C dx_i}{(df : f \in I)}.
  \]
  If $I = (f_1, \dots, f_r)$, then we just get
  \[
    \Omega_{C / k}
    = \frac{\bigoplus_{i = 1}^n C dx_i}{(df_1, \dots, df_r)}.
  \]
\end{example}

\begin{example}
  If $X = V(x^2 - y^3)$, then
  $C = \OO_X(X) = k[x, y] / (x^2 - y^3)$,
  and Example \ref{example:affine-variety-differentials} gives
  \[
    \Omega_{C / k}
    = \frac{Cdx \oplus C dy}{(2xdx - 3y^2 dy)}.
  \]
\end{example}

\begin{prop}
  Given a variety $X$, there exists a
  coherent sheaf $\Omega_X$,
  called the \emph{sheaf of differentials},
  such that
  \begin{enumerate}
    \item for $U \subseteq X$ open
      affine, there exists an isomorphism
      $\Omega_X(U) \cong \Omega_{\OO_X(U) / k}$,
    \item for $V \subseteq U \subseteq X$
      open affine, the following diagram
      commutes:
      \[
        \begin{tikzcd}
          \Omega_X(U) \ar[r] \ar[d] & \Omega_{\OO_X(U) / k} \ar[d, "df \mapsto d(f|_U)"] \\
          \Omega_X(V) \ar[r] & \Omega_{\OO_X(V) / k}
        \end{tikzcd}
      \]
  \end{enumerate}
\end{prop}

\begin{proof}
  For $U \subseteq X$ open affine,
  set $\mathcal{G}_U := \widetilde{\Omega_{\OO_X(U) / k}} \in \Coh(U)$.
  Now for $V \subseteq U \subseteq X$
  open affine,
  \[
    \Omega_{\OO_X(U) / k} \longrightarrow \Omega_{\OO_X(V) / k}
  \]
  induces a morphism
  $\phi_{V, U} : \mathcal{G}_U|_V
    \to \mathcal{G}_V$. If $V = D_U(f)$,
    then
    \[
      \mathcal{G}_U|_V
      = ((\Omega_{\OO_X(U) / k})_f)^\sim
      \quad \text{and} \quad
      \mathcal{G}_V
      = \widetilde{\Omega_{\OO_X(V) / k}}.
    \]
    These are isomorphic by the
    localization property of differentials.
    So by the affine communication
    lemma, we get that $\phi_{V, U}$ is
    an isomorphism for all $V \subseteq U \subseteq X$ open affine.

    Using this, we can glue the
    $\mathcal{G}_U$ for $U \subseteq X$
    open affine to construct
    $\mathcal{G} =: \Omega_{X}$.
\end{proof}

\begin{remark}
  Similarly, for a morphism of varieties
  $f : X \to Y$, we can define
  \[
    \Omega_{X / Y} \in \Coh(X)
  \]
  such that for any open affines
  $U \subseteq X$ and $V \subseteq Y$
  with $f(U) \subseteq V$, we have
  \[
    \Omega_{X / Y}|_U
    \cong (\Omega_{\OO_X(U) / \OO_Y(V)})^\sim.
  \]
\end{remark}

\begin{remark}
  Note that $\Omega_X \cong \Omega_{X / \{\mathrm{pt}\}}$.
\end{remark}

\begin{example}
  We have the following:
  \begin{enumerate}
    \item $\Omega_{\Affine^n} \cong (\Omega_{k[x_1, \dots, x_n] / k})^\sim \cong (\bigoplus_{i = 1}^n k[x_1, \dots, x_n] dx_i)^\sim \cong \bigoplus_{i = 1}^n \OO_{\Affine^n}$.
    \item To compute $\Omega_{\PP^n}$, note
      that there is a short exact sequence
      (called the \emph{Euler sequence})
      \[
        \begin{tikzcd}
          0 \ar[r] & \Omega_{\PP^n} \ar[r] & \OO_{\PP^n}(-1)^{\oplus (n + 1)} \ar[r] & \OO_{\PP^n} \ar[r] & 0
        \end{tikzcd}
      \]
      On $U_i \subseteq \PP^n$, this
      corresponds to
      \[
        \begin{tikzcd}[row sep=tiny]
          \Omega_{\PP^n}(U_i)
          \ar[r] & \bigoplus_{i = 1}^{n + 1} \OO_{\PP^n}(-1)(U_i)
          \ar[r] & \OO_{\PP^n}(U_i) \ar[r] & 0 \\
          d(x_j / z_i)
          \ar[r, mapsto]
                 & e_j - (z_j / z_i) e_i \\
                 & e_j \ar[r, mapsto] & z_j / z_i.
        \end{tikzcd}
      \]
      One can check that these maps glue
      and that the sequence is exact.
  \end{enumerate}
\end{example}
